% Standalone section draft; uses amsmath/amssymb/amsthm.
% Definitions: theorem, lemma, proposition, corollary environments in parent.
\section{The shortest positive average cancelling two frequencies}
\label{sec:two-frequency}

For a compactly supported Borel probability measure $\mu$ on $\mathbb R$, write
\[
 \phi_\mu(\omega)=\int_{\mathbb R}e^{i\omega t}\,d\mu(t),\qquad
 D(\mu)=\sup\operatorname{supp}\mu-\inf\operatorname{supp}\mu.
\]
Given $0<\tau<\sigma$, the problem is to minimize $D(\mu)$ subject to
$\phi_\mu(\tau)=\phi_\mu(\sigma)=0$. Translation multiplies each condition by a
nonzero phase, and scaling reduces the problem to $(\tau,\sigma)=(1,r)$,
where $r=\sigma/\tau>1$. Let $H(r)$ denote half the minimum diameter in this
normalization. The theorem below optimizes over \emph{all} probability measures;
neither finite support nor symmetry is an assumption.

\begin{theorem}[Complete two-frequency law]\label{thm:two-frequency}
The minimum diameter for $(\tau,\sigma)$ is $2H(r)/\tau$. Up to translation,
the minimizing measure is unique. In normalized coordinates it is as follows.
\begin{enumerate}
\item If $r$ is an odd integer, then $H(r)=\pi/2$ and
\[
 \mu_r=\tfrac12\delta_{-\pi/2}+\tfrac12\delta_{\pi/2}.
\]
\item Suppose $2k-1<r<2k+1$, where $k\ge1$ is an integer, and put
$s=(-1)^k$. If $k=1$, set $a=0$ and $M=1$. If $k\ge2$, let $a$ be the
unique solution of
\begin{equation}\label{eq:tangent}
 r\tan(ra)=\tan a,
 \qquad \frac{(k-1)\pi}{r}<a<\frac{(k-\tfrac12)\pi}{r},
\end{equation}
and set $M=-s\cos(ra)/\cos a$. There is a unique $H$ satisfying
\begin{equation}\label{eq:endpoint}
 \frac{s\cos(rH)}{-\cos H}=M,
 \qquad \frac\pi2<H<\min\left\{\pi,\frac{(k+\tfrac12)\pi}{r}\right\}.
\end{equation}
This $H$ is $H(r)$. Put
\[
 \alpha=\frac{-\cos H}{\cos a-\cos H},\qquad
 \beta=\frac{\cos a}{\cos a-\cos H}.
\]
Then $\alpha,\beta>0$, $\alpha+\beta=1$, and
\begin{equation}\label{eq:optimal-measure}
 \mu_r=\frac\alpha2(\delta_{-a}+\delta_a)
          +\frac\beta2(\delta_{-H}+\delta_H).
\end{equation}
When $a=0$, the first two masses combine to one atom.
\end{enumerate}
In particular, for $1<r<3$,
\begin{equation}\label{eq:low-ratio}
 H(r)=\frac{2\pi}{1+r},\qquad
 c=\cos\frac{2\pi}{1+r},\qquad
 \mu_r=\frac{\delta_{-H}+\delta_H}{2(1-c)}
             -\frac{c}{1-c}\delta_0.
\end{equation}
The same formula at $r=3$ reduces to the two endpoint atoms.
\end{theorem}

\begin{proof}
First, a measure on an interval shorter than $\pi$ cannot have
$\phi_\mu(1)=0$: after centering the interval, $\cos t>0$ throughout it.
At diameter $\pi$, positivity of $\cos t$ in $(-\pi/2,\pi/2)$ forces
support at the two endpoints, and the imaginary part forces equal masses.
Its second Fourier value is $\cos(r\pi/2)$, zero exactly when $r$ is odd.
This proves the first case and the strict lower bound $H(r)>\pi/2$ otherwise.

For the second case define
\[
 f(t)=\frac{-s\cos(rt)}{\cos t},\quad 0\le t<\pi/2.
\]
Its positive lobes are precisely the cosine lobes whose sign is $-s$.
Translating any earlier such lobe by an integer multiple of $2\pi/r$
to the last one preserves its numerator and decreases its positive
denominator. Within the last lobe, reflection about its centre increases
$t$, preserves the numerator, and again decreases the denominator.
Consequently its global maximum occurs on
\[
 I=\left[\frac{(k-1)\pi}{r},\frac{(k-\tfrac12)\pi}{r}\right).
\]
The maximum is positive; elsewhere, nonpositive values cannot compete.

For $k\ge2$, the logarithmic derivative of $f$ on $I$ is
$\tan t-r\tan(rt)$. Set $F(t)=r\tan(rt)-\tan t$.
At the left endpoint $F<0$ and at the right endpoint $F\to+\infty$.
At every zero of $F$ one has
\[
 F'(t)=r^2\sec^2(rt)-\sec^2t=r^2-1>0.
\]
Thus every zero is a strict crossing from negative to positive, and there
can be only one. This proves \eqref{eq:tangent}, that $a$ is the unique
maximum point, and $M=\max f$. For $k=1$ the same calculation, starting
at $F(0)=0$ and $F'(0)=r^2-1>0$, gives $F>0$ subsequently, so the unique
maximum is $a=0$ with $M=1$.

On
\[
 J=\left(\frac\pi2,\min\left\{\pi,
                    \frac{(k+\tfrac12)\pi}{r}\right\}\right)
\]
the function $R(t)=s\cos(rt)/(-\cos t)$ is positive. Its logarithmic
derivative is $-F(t)$, with the same expression for $F$. There are no
poles of $\tan(rt)$ inside $J$, and $F(t)\to+\infty$ as
$t\downarrow\pi/2$. The identity $F'=r^2-1>0$ at every zero again
prevents a crossing from positive to negative. Hence $R$ is strictly
decreasing. It starts at $+\infty$. It tends to zero at the right
endpoint unless that endpoint is $\pi$; the latter case occurs only
for $1<r\le3/2$, and then $R(\pi)=-\cos(r\pi)<1=M$.
The endpoint equation therefore has exactly one solution.

The decisive dual function is
\begin{equation}\label{eq:dual}
 q(t)=s\cos(rt)+M\cos t.
\end{equation}
For $0\le t<\pi/2$, maximality of $M$ gives $q(t)\ge0$, with equality
only at $a$. For $\pi/2<t\le H$, the inequality $R(t)\ge R(H)=M$
gives $q(t)\ge0$, with equality only at $H$. At $\pi/2$ it is positive.
By evenness, $q\ge0$ on $[-H,H]$ and its zero set there is exactly
$\{-H,-a,a,H\}$. The identities
$\cos(ra)=-sM\cos a$ and $\cos(rH)=-sM\cos H$ show directly that
\eqref{eq:optimal-measure} has both required Fourier values zero.

Now centre the support of any competitor of diameter at most $2H$.
The constraints give $\int q\,d\mu=0$, so positivity forces its support
onto the displayed zero set. Its total masses on $\{-a,a\}$ and
$\{-H,H\}$ are necessarily $\alpha$ and $\beta$, by mass one and the
real first-frequency constraint. If $a>0$, write $d_a,d_H$ for the
differences between positive and negative masses at the corresponding
locations. The imaginary constraints are
\[
 d_a\sin a+d_H\sin H=0,\qquad
 d_a\sin(ra)+d_H\sin(rH)=0.
\]
Since $q'(a)=0$, their determinant is
\[
 \sin a\sin(rH)-\sin H\sin(ra)
   =-\frac{s\sin a}{r}q'(H)\ne0.
\]
Here $q'(H)<0$ because $R'(H)<0$. Thus $d_a=d_H=0$.
If $a=0$, the first imaginary constraint already makes the endpoint
masses equal. This proves uniqueness. Both endpoint masses are
positive, so a competitor with smaller diameter is impossible.

Finally, if $1<r<3$, the equation $\cos(rH)=\cos H$ has its required
solution at $H=2\pi/(1+r)$, giving \eqref{eq:low-ratio}.
\end{proof}

\paragraph{An exact four-atom example.}
At $r=4$, put
\[
 x=\sqrt{\frac{2+\sqrt{10}}{12}},\qquad
 z=x-\sqrt{1-2x^2}.
\]
Then the theorem gives
\[
 a=\arccos x,\qquad H=\arccos(-z),\qquad M=16x(1-2x^2).
\]
Indeed, $24x^4-8x^2-1=0$ is the tangent equation, and
\[
 8X^4-8X^2+MX+1
   =8(X-x)^2\{X^2+2xX+(3x^2-1)\}
\]
identifies the endpoint root $X=-z$ exactly. For orientation,
$a\approx0.8554371830$, $H\approx1.8569146033$, with total interior
and exterior masses approximately $0.3008471646$ and $0.6991528354$.
The radical formulas, not these decimals, specify the example.

\subsection{Resonance produces a cube-root cusp}

The simple-looking exact resonances have a singular neighbourhood.
Let $m\ge3$ be odd and define
\begin{equation}\label{eq:resonance-constants}
 A_m=m(m^2-1),\qquad
 C_m=\left(\frac{3\pi}{2A_m}\right)^{1/3},\qquad
 M_m^- =\max_{0<\delta\le\pi/2}
             \frac{-\sin(m\delta)}{\sin\delta}.
\end{equation}
The last maximum is finite and satisfies $0<M_m^-<m$. For example,
$M_3^-=1$ and $M_5^-=5/4$.

\begin{theorem}[Asymmetric resonance law]\label{thm:cusp}
For every odd $m\ge3$, as $\varepsilon\downarrow0$,
\begin{align}
 H(m-\varepsilon)
   &=\frac\pi2+\frac{\pi}{2(m+M_m^-)}\varepsilon
                       +o(\varepsilon),\label{eq:left-cusp}\\
 H(m+\varepsilon)
   &=\frac\pi2+2C_m\varepsilon^{1/3}
                       +o(\varepsilon^{1/3}).\label{eq:right-cusp}
\end{align}
On the right, the interior node and dual coefficient satisfy
\[
 a(m+\varepsilon)=\frac\pi2-C_m\varepsilon^{1/3}
                                  +o(\varepsilon^{1/3}),\qquad
 M(m+\varepsilon)=m-\frac{A_mC_m^2}{2}\varepsilon^{2/3}
                                  +o(\varepsilon^{2/3}),
\]
and the total masses of the interior and exterior pairs tend to $2/3$
and $1/3$, respectively. In particular, the optimal support radius is
not locally Lipschitz at any such odd resonance, and no pointwise
H\"older estimate with exponent greater than $1/3$ can hold there.
\end{theorem}

\begin{proof}
Write $r=m+\varepsilon$, $\eta=\pi\varepsilon/2$,
$u=\pi/2-a$ and $v=H-\pi/2$. On this branch the tangent and endpoint
equations become
\begin{align}
 r\cos(ru-\eta)\sin u-\sin(ru-\eta)\cos u&=0,
                                                       \label{eq:cusp-tangent}\\
 \frac{\sin(\eta+rv)}{\sin v}=M
     =\frac{\sin(ru-\eta)}{\sin u}.       \label{eq:cusp-endpoint}
\end{align}
The allowed interval for $u$ has limit $[0,\pi/(2m)]$.
A positive limit $u_*$ is impossible in \eqref{eq:cusp-tangent}:
for $0<u<\pi/(2m)$, strict increase of $\tan x/x$ gives
$\tan(mu)>m\tan u$, and the endpoint $\pi/(2m)$ is also not a root.
Thus $u\to0$.

Taylor expansion of \eqref{eq:cusp-tangent}, with a remainder uniform
for small $\eta,u$, gives
\[
 0=\eta-\frac{A_m}{3}u^3+o(\eta+u^3).
\]
The two nonnegative terms must balance, so
$\eta\sim A_mu^3/3$ and $u\sim C_m\varepsilon^{1/3}$.
The tangent equation also gives the useful identity
\[
 M=\frac{r\cos(ru-\eta)}{\cos u}
   =r-\frac{A_m}{2}u^2+o(u^2).
\]
Hence $M\to m$. Any positive subsequential limit $v_*$ in the allowed
endpoint interval would satisfy
$\sin(mv_*)/\sin v_*=m$, contrary to
$|\sin(mv)|<m\sin v$ for $0<v<\pi$ where the numerator is positive.
Consequently $v\to0$.

Expanding the endpoint equation without division by $v$ yields
\[
 0=\eta+\left(\frac{A_m}{2}+o(1)\right)u^2v
            -\left(\frac{A_m}{6}+o(1)\right)v^3
            +o(\eta).
\]
This equation and $\eta\sim A_mu^3/3$ first imply that $v/u$ is bounded
above and bounded away from zero: otherwise divide respectively by
$v^3$ or by $u^3$. Every limit $z>0$ then satisfies
$z^3-3z-2=(z-2)(z+1)^2=0$. Thus $v/u\to2$.
Finally,
$\alpha=\sin v/(\sin u+\sin v)\to2/3$ and
$\beta=\sin u/(\sin u+\sin v)\to1/3$.

For the left side put $r=m-\varepsilon$ and again $\eta=\pi\varepsilon/2$.
The interval containing the maximizing $a$ stays a positive distance
from $\pi/2$, and continuous dependence of its unique maximum gives
$M(r)\to M_m^-$. The endpoint lies below $m\pi/(2r)$, hence
$v=O(\varepsilon)$. Its equation is now
\[
 \sin(\eta-rv)=M(r)\sin v.
\]
It follows that $\eta=(r+M(r))v+O(\varepsilon^3)$, proving
\eqref{eq:left-cusp}. Positivity of $M_m^-$ follows by taking
$\delta=3\pi/(2m)$, and its strict upper bound follows from the strict
sine inequality away from zero and its negative limit $-m$ at zero.
\end{proof}

\begin{corollary}[Merging distinct frequencies has a discontinuous cost]
\label{cor:merger}
If the repeated frequency is imposed only once, $H(1)=\pi/2$.
Nevertheless, $\lim_{r\downarrow1}H(r)=\pi$. The optimal measures for
distinct coalescing frequencies converge weakly to
\[
 \tfrac14\delta_{-\pi}+\tfrac12\delta_0+\tfrac14\delta_{\pi},
\]
whose Fourier transform is $\cos^2(\pi\omega/2)$ and therefore has a
double zero at $\omega=1$.
\end{corollary}

\begin{proof}
All assertions follow directly from \eqref{eq:low-ratio}. The limiting
measure is also the unique shortest positive measure with
$\phi(1)=\phi'(1)=0$, by the one-frequency double-zero result in the
source article (Theorem \texttt{rec:thm:fwidth}).
\end{proof}

\subsection{Stability, detuning, and common real parts}

\begin{theorem}[Exact first-order price of shortening]
\label{thm:shortening}
For $r>1$ and $0<d<H(r)$ set
\[
 D_r(d)=\inf_{\operatorname{supp}\mu\subset[-H(r)+d,H(r)-d]}
               \max\{|\phi_\mu(1)|,|\phi_\mu(r)|\},
\]
where the infimum runs over probability measures. If $r$ is not odd,
let $a,H,M,s,\beta$ be as in Theorem \ref{thm:two-frequency} and
$q(t)=M\cos t+s\cos(rt)$. Then
\begin{equation}\label{eq:shortening-slope}
 D_r(d)=\kappa_r d+O_r(d^2),\qquad
 \kappa_r=\frac{-\beta q'(H)}{M+1}>0.
\end{equation}
If $r=m\ge3$ is odd, then instead
\begin{equation}\label{eq:odd-shortening}
 D_m(d)=\frac{m+M_m^-}{1+M_m^-}d+O_m(d^2).
\end{equation}
Thus $D_3(d)=2d+O(d^2)$ and $D_5(d)=\tfrac{25}{9}d+O(d^2)$.
The parameter $d$ is the loss in \emph{half-width}; the diameter is
shortened by $2d$.
\end{theorem}

\begin{proof}
For the nonodd case put $K=-\beta q'(H)>0$ and consider
\[
 q_d(t)=q(t)+d\frac{K}{\cos a}\cos t.
\]
The interior zeros of $q$ are quadratic, with
$q''(a)=(r^2-1)M\cos a>0$. Its endpoint zeros are simple. At each
interior contact the perturbation has value $dK$. At the truncated
right endpoint,
\[
 q_d(H-d)
   =d\left(-q'(H)+\frac{K\cos H}{\cos a}\right)+O_r(d^2)
   =dK+O_r(d^2),
\]
because $\beta=\cos a/(\cos a-\cos H)$. The same holds on the left.
Taylor's theorem near the finitely many contacts and a positive lower
bound away from them therefore give
$\min_{|t|\le H-d}q_d(t)\ge dK-O_r(d^2)$.
For a measure on that interval with residual $e$,
\[
 dK-O_r(d^2)\le\int q_d\,d\mu
       \le\left(M+1+d\frac{K}{\cos a}\right)e.
\]
This proves the lower bound in \eqref{eq:shortening-slope}.

For the matching upper bound, keep the interior nodes at $\pm a$,
move the exterior nodes to $\pm(H-d)$, and change the total exterior
mass to $\beta+bd$. Set $\Delta=\cos H-\cos a<0$ and choose
\[
 b\Delta=\frac{\beta\{sr\sin(rH)-\sin H\}}{M+1}.
\]
All masses remain positive for small $d$. Their two real Fourier
values have the first-order expansions
\[
 d\{\beta\sin H+b\Delta\}+O_r(d^2),\qquad
 d\{\beta r\sin(rH)-sMb\Delta\}+O_r(d^2).
\]
The coefficients are respectively $\kappa_r$ and $s\kappa_r$.
The imaginary parts vanish, completing the nonodd case.

For an odd $m$, let $a_*$ maximize
$-\sin(m\pi/2)\cos(mt)/\cos t$ on $0\le t<\pi/2$.
Its value is $M=M_m^-$, and $a_*<\pi/2$; the preceding lobe argument
proves uniqueness. The function
\[
 q(t)=M\cos t+\sin(m\pi/2)\cos(mt)
\]
is nonnegative on $[-\pi/2,\pi/2]$, has quadratic interior zeros at
$\pm a_*$, and has endpoint derivative $q'(\pi/2)=-(M+m)$.
The same lower-bound perturbation, with $K=M+m$ and $a=a_*$, gives
$D_m(d)\ge (M+m)d/(M+1)-O_m(d^2)$.
For the upper bound, put total mass
\[
 \alpha_d=\frac{m-1}{(M+1)\cos a_*}\,d
\]
at the symmetric interior pair and mass $1-\alpha_d$ at the truncated
endpoint pair. The two real moments are
$\frac{M+m}{M+1}d+O_m(d^2)$ and
$\sin(m\pi/2)\frac{M+m}{M+1}d+O_m(d^2)$, respectively.
This proves \eqref{eq:odd-shortening}.
\end{proof}

\begin{corollary}[A second singularity at resonance]
For odd $m\ge3$ the shortening coefficient satisfies
\[
 \kappa_{m-\varepsilon}\longrightarrow
       \frac{m+M_m^-}{1+M_m^-}=\kappa_m,
 \qquad
 \kappa_{m+\varepsilon}\sim
       \frac{m(m-1)}2C_m^2\varepsilon^{2/3}.
\]
In particular, the limits $d\downarrow0$ and $r\downarrow m$ do not
commute for $D_r(d)/d$.
\end{corollary}

\begin{proof}
The left limit follows from the continuity used in
\eqref{eq:left-cusp}. On the right, the notation in the proof of
Theorem \ref{thm:cusp} gives
$q'(H)=-3A_mu^2/2+o(u^2)$ and $\beta\to1/3$.
Substitute in \eqref{eq:shortening-slope} and use
$u\sim C_m\varepsilon^{1/3}$.
For fixed $0<d<\pi/2$, compactness of probability measures on a common
compact interval, together with linear rescaling of the support
intervals, gives $D_r(d)\to D_m(d)$ as $r\to m$.
This justifies the assertion about the two iterated limits.
\end{proof}

\begin{proposition}[Sharp stability at the minimum support]
\label{prop:filter-stability}
Fix a nonodd $r>1$ and let $H=H(r)$. There is a constant $C_r$ such that
every probability measure supported on $[-H,H]$ satisfies
\[
 W_1(\mu,\mu_r)\le C_r
       \max\{|\phi_\mu(1)|,|\phi_\mu(r)|\}^{1/2}.
\]
The exponent $1/2$ is optimal. At an odd integer $r$, the same statement
holds with exponent $1$, which is optimal.
\end{proposition}

\begin{proof}
For nonodd $r$, the dual function \eqref{eq:dual} has simple zeros at
the endpoints and quadratic zeros at the interior nodes. Indeed,
\[
 q''(a)=(r^2-1)M\cos a>0
\]
also when $a=0$, and $q'(H)<0$. Thus
$q(t)\ge c_r\operatorname{dist}(t,G)^2$ on $[-H,H]$, where
$G=\{-H,-a,a,H\}$. If the displayed Fourier residual is $e$, then
$0\le\int q\,d\mu\le(M+1)e$. Rounding every point to a closest
node costs $O_r(\sqrt e)$ in $W_1$ and changes both Fourier moments
by $O_r(\sqrt e)$. The map from signed masses on $G$ to total mass and
the two complex Fourier values is injective, by the real and imaginary
calculations in the uniqueness proof. On this finite-dimensional
space its inverse on the image is bounded. The rounded masses are
therefore $O_r(\sqrt e)$ from those of $\mu_r$.
For sharpness, split one positive interior mass equally between
points at distance $d$ on either side. The residual is $O_r(d^2)$,
whereas its $W_1$ distance from $\mu_r$ is bounded below by a positive
multiple of $d$.

At an odd ratio, use $q(t)=\cos t$, which is bounded below by a
positive constant times distance to $\{-\pi/2,\pi/2\}$, and round
to these two endpoints. Their imbalance is controlled by the first
imaginary moment. Moving both endpoints inward by $d$ gives residual
and transport distance of order $d$, proving optimality.
\end{proof}

\begin{proposition}[Quadratic suppression between the cancelled frequencies]
\label{prop:detuning}
If an even probability measure is supported on $[-L,L]$ and its
Fourier transform vanishes at two distinct positive frequencies
$0<\tau<\sigma$, then for every real $\omega$,
\[
 |\phi_\mu(\omega)|\le
 \min\left\{1,\frac{L^2}{2}|\omega-\tau||\omega-\sigma|\right\}.
\]
In particular its modulus is at most $L^2(\sigma-\tau)^2/8$ throughout
the interval between the two prescribed frequencies. Using evenness
and all four zeros also gives
\[
 |\phi_\mu(\omega)|\le
 \min\left\{1,\frac{L^4}{24}
       |(\omega^2-\tau^2)(\omega^2-\sigma^2)|\right\}.
\]
\end{proposition}

\begin{proof}
The Fourier transform is real on the real axis and its $j$th derivative
has modulus at most $L^j$. Apply the real interpolation remainder
formula to the degree-one zero interpolant at $\tau,\sigma$, or to
the degree-three zero interpolant at $\pm\tau,\pm\sigma$.
The assertion remains true at an interpolation node by continuity.
\end{proof}

\begin{corollary}[Cancellation on any common vertical line]
\label{cor:tilting}
Fix $0<\tau<\sigma$ and $b\in\mathbb R$, and define
$F_\nu(z)=\int e^{zt}\,d\nu(t)$. The least diameter of a probability
measure for which $F_\nu(b+i\tau)=F_\nu(b+i\sigma)=0$ is still
$2H(\sigma/\tau)/\tau$. Its unique minimizer up to translation is the
exponential tilt of the measure in Theorem \ref{thm:two-frequency}:
if the normalized optimal atoms are $t_j$ with masses $w_j$, use
locations $t_j/\tau$ and masses
\[
 \frac{w_j e^{-bt_j/\tau}}{\sum_\ell w_\ell e^{-bt_\ell/\tau}}.
\]
\end{corollary}

\begin{proof}
The map $d\mu=e^{bt}\,d\nu/\int e^{bu}\,d\nu(u)$ is a bijection of
probability measures with the same support and converts the two
Laplace-transform constraints to the two Fourier constraints.
\end{proof}

\subsection{Application to the normalized prime errors in the article}

The input GRH theorem is used here only to put the relevant zeros on
one vertical line. Fix a modulus $q$, a reduced residue class $b$, and
write the source article's ramp statistic as
\[
 S_b(x)=\sum_{\substack{n\le x\\n\equiv b\pmod q}}
              \Lambda(n)\left(\frac1n-\frac1x\right).
\]
With the constants $c_b$ and $\beta_b$ of its
\texttt{zer:eq:va}, define
\[
 Z_b(t)=e^{t/2}\left(S_b(e^t)-\frac{t}{\varphi(q)}-c_b\right)
                     -\beta_b.
\]
Its Lemma \texttt{zer:lem:ramp}, equation \texttt{zer:eq:Zexp}, gives
under Dirichlet GRH
\[
 Z_b(t)=\sum_{\substack{\gamma\in\Gamma_q\\\gamma\ne0}}
           \frac{v_b(\gamma)}{1/4+\gamma^2}e^{i\gamma t}
             +O_q((1+t)e^{-t/2}),
\]
where $\Gamma_q$ groups equal ordinates, multiplicities are included
in $v_b$, and the series is absolutely and uniformly convergent.

\begin{corollary}[Shortest positive two-frequency deletion]
\label{cor:prime-deletion}
Take $0<\tau<\sigma$ with $\tau,\sigma\in\Gamma_q$ and set
$L=H(\sigma/\tau)/\tau$. Scale the optimal measure of
Theorem \ref{thm:two-frequency} to $[-L,L]$ and call it $\nu$.
For $t\ge L+\log2$, simultaneously for every reduced residue class,
\[
 \int Z_b(t+u)\,d\nu(u)
 =\sum_{\substack{\gamma\in\Gamma_q\\
         \gamma\notin\{0,\pm\tau,\pm\sigma\}}}
       \frac{v_b(\gamma)\phi_\nu(\gamma)}{1/4+\gamma^2}
                     e^{i\gamma t}
       +O_q\left((1+t+L)e^{-(t-L)/2}\right).
\]
The averaging rule uses at most four scales, all in
$[xe^{-L},xe^L]$ when $x=e^t$, and its ratio of largest to smallest
scale is $e^{2L}$. No positive translation-independent rule with
smaller support diameter has Fourier multiplier zero at both
frequencies. No simplicity, separation from other ordinates, or
rational independence of zeros is required.
\end{corollary}

\begin{proof}
Absolute convergence permits termwise integration and the multiplier
at a frequency $\gamma$ is $\phi_\nu(\gamma)$. The four prescribed
frequencies vanish. The remainder follows by its uniform estimate
over $|u|\le L$. Optimality is exactly Theorem
\ref{thm:two-frequency}.
\end{proof}

This assertion concerns optimal deletion of two specified spectral
modes. It does not assert a new bound for the unfiltered arithmetic
remainder, certify any numerically computed zero, or cancel the
remaining spectrum. If instead one averages an unnormalized term
$x^{b+i\gamma}$, Corollary \ref{cor:tilting} supplies the positive
weights: use $b=-1/2$ for the source ramp remainder, or $b=1/2$ for
the usual unnormalized explicit-formula zero terms.

\paragraph{Relation to the source article and prior work.}
The article's Theorem \texttt{rec:thm:fwidth} treats one frequency
with several prescribed real parts and multiplicities. Its Theorem
\texttt{rec:thm:hblock} treats consecutive harmonic blocks. The
question at the end of \texttt{rec:sec:open} asks about nonconsecutive
prescribed frequencies. The present result gives the complete answer
for two distinct positive real frequencies, including irrational ratios,
with a singular law at every odd integer ratio $m\ge3$.
The three-atom construction for $1<r\le3$ is already present in
classical two-mode input shaping: see Convolve, Inc.,
\emph{Minimizing Structural Vibrations with Input Shaping},
NASA-CR-189439 (January 1995), printed pages 39--42, especially
equations (15)--(16) \cite{NASA1995}; and W. Singhose, E. Crain, and W. Seering,
\emph{Convolved and simultaneous two-mode input shapers},
IEE Proceedings--Control Theory and Applications 144 (1997),
doi:10.1049/ip-cta:19971439 \cite{Shaping1997}. This low-ratio formula must not be
advertised as new. The proposed contribution is the all-ratio
positive-measure extremal theorem, its full equality classification,
and the one-sided cube-root resonance asymptotics. A search is
evidence about priority, not a proof of priority.
